% !TeX root = propuesta_prueba_semana04.tex
% !TeX program = pdflatex
% Propuesta parcial de la primera evaluación.
\documentclass[10pt,a5paper]{article}
\newif\ifmostrarSoluciones
\mostrarSolucionestrue
\usepackage{estilo_apuntes}
\encabezado{Semana 4}{Propuesta de evaluación}

\begin{document}
\tituloclase{Primera evaluación}{Propuesta parcial de ejercicios}{Semanas 1 a 3}

\noindent
\textbf{Nombre:} \rule{7cm}{0.4pt}

\medskip
\noindent
\textbf{Fecha:} \rule{3cm}{0.4pt}
\hfill
\textbf{Puntaje:} \rule{1.5cm}{0.4pt}

\begin{nota}[indicaciones]
Justifique cada procedimiento. Una respuesta sin el análisis de los límites laterales, las restricciones algebraicas o las cotas utilizadas recibirá únicamente crédito parcial. Puntaje sugerido: 10 puntos por ejercicio.
\end{nota}

\begin{ejercicios}[propuesta parcial]
\begin{enumerate}
  \item Sean $a,b\in\mathbb{R}$ y sea $F\colon\mathbb{R}\to\mathbb{R}$ la función definida por
  \[
    F(x)=
    \begin{cases}
      \dfrac{x^2-1}{x-1}, & x<1,\\[2mm]
      5, & x=1,\\
      ax+b, & x>1.
    \end{cases}
  \]
  Determine $a$ y $b$ para que exista
  \[
    \lim_{x\to1}F(x)
  \]
  y, además, se cumpla $F(2)=5$. Calcule el límite y compárelo con $F(1)$.

  \item Sea $B$ una función definida cerca de $0$, excepto posiblemente en $0$, que satisface
  \[
    |B(x)|\leq4.
  \]
  Calcule mediante el teorema del Sanduche
  \[
    \lim_{x\to0}
    \frac{\sqrt{1+x^2}-1}{x}\,B(x).
  \]
  No suponga que $B$ tiene límite.

  \item Considere las funciones
  \[
    f\colon[1,\infty)\to[0,\infty),
    \qquad
    f(x)=\sqrt{x-1},
  \]
  y
  \[
    g\colon\mathbb{R}\to[1,\infty),
    \qquad
    g(x)=x^2+1.
  \]
  \begin{enumerate}
    \item Determine si cada función es inyectiva, sobreyectiva o biyectiva.
    \item Halle la ley de asignación y el dominio de $g\circ f$ y de $f\circ g$.
    \item Decida si $f$ y $g$ son funciones inversas. Justifique su respuesta.
  \end{enumerate}
\end{enumerate}
\end{ejercicios}

\ifmostrarSoluciones
\begin{solucion}[de la propuesta]
\begin{enumerate}
  \item Para $x<1$ y $x\neq1$,
  \[
    \frac{x^2-1}{x-1}=x+1.
  \]
  Por ello,
  \[
    \lim_{x\to1^-}F(x)=2.
  \]
  Por la derecha,
  \[
    \lim_{x\to1^+}F(x)=a+b.
  \]
  La existencia del límite exige
  \[
    a+b=2.
  \]
  Además, la condición $F(2)=5$ implica
  \[
    2a+b=5.
  \]
  Al resolver el sistema se obtiene
  \[
    a=3,
    \qquad
    b=-1.
  \]
  En consecuencia,
  \[
    \lim_{x\to1}F(x)=2,
    \qquad
    F(1)=5.
  \]
  El límite y el valor de la función en el punto son diferentes.

  \item Para $x\neq0$, racionalizamos el factor algebraico:
  \[
    \frac{\sqrt{1+x^2}-1}{x}
    =\frac{x}{\sqrt{1+x^2}+1}.
  \]
  Entonces
  \begin{align*}
    \left|
      \frac{\sqrt{1+x^2}-1}{x}B(x)
    \right|
    &=\frac{|x|}{\sqrt{1+x^2}+1}|B(x)|\\
    &\leq\frac{4|x|}{\sqrt{1+x^2}+1}\\
    &\leq2|x|.
  \end{align*}
  Por tanto,
  \[
    -2|x|
    \leq
    \frac{\sqrt{1+x^2}-1}{x}B(x)
    \leq
    2|x|.
  \]
  Como ambas cotas tienden a $0$, el teorema del Sanduche permite concluir que
  \[
    \lim_{x\to0}
    \frac{\sqrt{1+x^2}-1}{x}B(x)=0.
  \]

  \item La función $f$ es estrictamente creciente en $[1,\infty)$, por lo que es
  inyectiva. Además, para cada $y\in[0,\infty)$, el punto
  \[
    x=y^2+1
  \]
  pertenece a $[1,\infty)$ y satisface $f(x)=y$. Por tanto, $f$ es sobreyectiva
  y, en consecuencia, biyectiva.

  La función $g$ es sobreyectiva sobre $[1,\infty)$: dado $y\geq1$, basta tomar
  \[
    x=\sqrt{y-1}.
  \]
  Sin embargo, no es inyectiva porque, por ejemplo,
  \[
    g(1)=g(-1)=2.
  \]
  Por ello, $g$ no es biyectiva.

  Como $f(x)\in[0,\infty)\subseteq\mathbb{R}$, la composición $g\circ f$ está
  definida para todo $x\in[1,\infty)$ y
  \[
    (g\circ f)(x)
    =\left(\sqrt{x-1}\right)^2+1
    =x.
  \]
  Por tanto,
  \[
    \operatorname{Dom}(g\circ f)=[1,\infty).
  \]

  Como $g(x)\in[1,\infty)$ para todo $x\in\mathbb{R}$, la composición
  $f\circ g$ está definida en $\mathbb{R}$ y
  \[
    (f\circ g)(x)
    =\sqrt{x^2}
    =|x|.
  \]
  Así,
  \[
    \operatorname{Dom}(f\circ g)=\mathbb{R}.
  \]
  Aunque $g\circ f$ es la identidad en $[1,\infty)$, $f\circ g$ no es la
  identidad en $\mathbb{R}$. Por tanto, $f$ y $g$ no son inversas con los
  dominios dados. Si se restringe el dominio de $g$ a $[0,\infty)$, entonces sí
  son funciones inversas.
\end{enumerate}
\end{solucion}
\fi

\end{document}
